
\documentclass{article}
\usepackage{colortbl}
\usepackage{makecell}
\usepackage{multirow}
\usepackage{supertabular}

\begin{document}

\newcounter{utterance}

\centering \large Interaction Transcript for game `dond', experiment `semi\_en', episode 15 with Qwen3{-}32B{-}t0.0.
\vspace{24pt}

{ \footnotesize  \setcounter{utterance}{1}
\setlength{\tabcolsep}{0pt}
\begin{supertabular}{c@{$\;$}|p{.15\linewidth}@{}p{.15\linewidth}p{.15\linewidth}p{.15\linewidth}p{.15\linewidth}p{.15\linewidth}}
    \# & $\;$A & \multicolumn{4}{c}{Game Master} & $\;\:$B\\
    \hline

    \theutterance \stepcounter{utterance}  
    & & \multicolumn{4}{p{0.6\linewidth}}{
        \cellcolor[rgb]{0.9,0.9,0.9}{
            \makecell[{{p{\linewidth}}}]{
                \texttt{\tiny{[P1$\langle$GM]}}
                \texttt{You are playing a negotiation game in which you have to agree on how to divide a set of items among you and another player.} \\
\\ 
\texttt{Rules:} \\
\texttt{(a) You and the other player are given a set of items. Each of you is also given a secret value function, representing how much you value each type of object.} \\
\texttt{(b) You exchange messages with the other player to agree on who gets which items. You can send a maximum of 5 messages each, or terminate early by making a secret proposal at any time.} \\
\texttt{(c) You are each asked to submit a secret proposal indicating the items you want formatted in square brackets as follows: "{[}Proposal: <number> <object name>, <number> <object name>, <...>{]}"} \\
\texttt{(d) If your proposals are complementary, i.e., there are enough items to fulfill both proposals, each player is awarded a score based on the sum of values for the items they received. Otherwise, both of you get zero points.} \\
\texttt{(e) Your goal is to maximize the score you receive.} \\
\\ 
\texttt{Let us start.} \\
\\ 
\texttt{The set of available items is:} \\
\\ 
\texttt{1 bowl, 1 bottle, 2 fruits, 1 eraser} \\
\\ 
\texttt{Your secret value function is:} \\
\\ 
\texttt{bowl: 5, bottle: 5, fruit: 0, eraser: 0} \\
\\ 
\texttt{IMPORTANT: Your messages, unless it is the secret proposal, are directly transmitted to the other player, so do not include any response to the rules or text announcing your message. To make a secret proposal, use the indicated format. Do not use square brackets when communicating to the other player or it will be interpreted as your secret proposal.} \\
            }
        }
    }
    & \\ \\

    \theutterance \stepcounter{utterance}  
    & \multicolumn{4}{p{0.6\linewidth}}{
        \cellcolor[rgb]{0.9,0.9,0.9}{
            \makecell[{{p{\linewidth}}}]{
                \texttt{\tiny{[P1$\rangle$GM]}}
                \texttt{<think>} \\
\texttt{Okay, let's see. I need to negotiate the division of these items: 1 bowl, 1 bottle, 2 fruits, 1 eraser. My value function is bowl:5, bottle:5, fruits and eraser:0. So I care a lot about getting the bowl and the bottle, but the fruits and eraser are worthless to me. The other player's value function is secret, so I don't know what they want. } \\
\\ 
\texttt{First, I should figure out what the other player might value. Since the items are bowl, bottle, fruits, and eraser, maybe they value the fruits and eraser? Or maybe they also want the bowl and bottle. But since there's only one of each, if both of us want them, we can't both get them. So the key is to find out what they want and see if we can split the items in a way that our proposals are complementary.} \\
\\ 
\texttt{The rules say that if our secret proposals are complementary, we get our scores. Otherwise, zero. So I need to make sure that whatever I propose, the other player's proposal doesn't conflict with mine. Since I don't know their value function, I have to make an educated guess.} \\
\\ 
\texttt{Let me think. The available items are limited. If I take the bowl and bottle, which are the only items I value, then the other player would have to take the fruits and eraser. But maybe they want some of the high{-}value items too. If they also want the bowl or bottle, then we might have a conflict. But since I can't get points from fruits or eraser, maybe I can offer them to the other player in exchange for the bowl and bottle. } \\
\\ 
\texttt{Wait, but how do I communicate that? I can send messages to negotiate. Let me start by sending a message to the other player. Maybe I can propose that I take the bowl and bottle, and they take the rest. But I need to phrase it in a way that encourages them to accept. } \\
\\ 
\texttt{So my first message could be: "I would like the bowl and the bottle. I don't value the fruits or eraser, so I'm happy for you to have them. Let me know if that works for you." That way, I'm stating my preference and offering the other items to them. If they don't need the bowl and bottle, they can take the fruits and eraser, and our proposals would be complementary} \\
            }
        }
    }
    & & \\ \\

    \theutterance \stepcounter{utterance}  
    & & & \multicolumn{2}{p{0.3\linewidth}}{
        \cellcolor[rgb]{0.9,0.9,0.9}{
            \makecell[{{p{\linewidth}}}]{
                \texttt{\tiny{[GM$|$GM]}}
                \texttt{continue} \\
            }
        }
    }
    & & \\ \\

    \theutterance \stepcounter{utterance}  
    & & & \multicolumn{2}{p{0.3\linewidth}}{
        \cellcolor[rgb]{0.9,0.9,0.9}{
            \makecell[{{p{\linewidth}}}]{
                \texttt{\tiny{[GM$|$GM]}}
                \texttt{<think>} \\
\texttt{Okay, let's see. I need to negotiate the division of these items: 1 bowl, 1 bottle, 2 fruits, 1 eraser. My value function is bowl:5, bottle:5, fruits and eraser:0. So I care a lot about getting the bowl and the bottle, but the fruits and eraser are worthless to me. The other player's value function is secret, so I don't know what they want. } \\
\\ 
\texttt{First, I should figure out what the other player might value. Since the items are bowl, bottle, fruits, and eraser, maybe they value the fruits and eraser? Or maybe they also want the bowl and bottle. But since there's only one of each, if both of us want them, we can't both get them. So the key is to find out what they want and see if we can split the items in a way that our proposals are complementary.} \\
\\ 
\texttt{The rules say that if our secret proposals are complementary, we get our scores. Otherwise, zero. So I need to make sure that whatever I propose, the other player's proposal doesn't conflict with mine. Since I don't know their value function, I have to make an educated guess.} \\
\\ 
\texttt{Let me think. The available items are limited. If I take the bowl and bottle, which are the only items I value, then the other player would have to take the fruits and eraser. But maybe they want some of the high{-}value items too. If they also want the bowl or bottle, then we might have a conflict. But since I can't get points from fruits or eraser, maybe I can offer them to the other player in exchange for the bowl and bottle. } \\
\\ 
\texttt{Wait, but how do I communicate that? I can send messages to negotiate. Let me start by sending a message to the other player. Maybe I can propose that I take the bowl and bottle, and they take the rest. But I need to phrase it in a way that encourages them to accept. } \\
\\ 
\texttt{So my first message could be: "I would like the bowl and the bottle. I don't value the fruits or eraser, so I'm happy for you to have them. Let me know if that works for you." That way, I'm stating my preference and offering the other items to them. If they don't need the bowl and bottle, they can take the fruits and eraser, and our proposals would be complementary} \\
            }
        }
    }
    & & \\ \\

    \theutterance \stepcounter{utterance}  
    & & \multicolumn{4}{p{0.6\linewidth}}{
        \cellcolor[rgb]{0.9,0.9,0.9}{
            \makecell[{{p{\linewidth}}}]{
                \texttt{\tiny{[GM$\rangle$P2]}}
                \texttt{You are playing a negotiation game in which you have to agree on how to divide a set of items among you and another player.} \\
\\ 
\texttt{Rules:} \\
\texttt{(a) You and the other player are given a set of items. Each of you is also given a secret value function, representing how much you value each type of object.} \\
\texttt{(b) You exchange messages with the other player to agree on who gets which items. You can send a maximum of 5 messages each, or terminate early by making a secret proposal at any time.} \\
\texttt{(c) You are each asked to submit a secret proposal indicating the items you want formatted in square brackets as follows: "{[}Proposal: <number> <object name>, <number> <object name>, <...>{]}"} \\
\texttt{(d) If your proposals are complementary, i.e., there are enough items to fulfill both proposals, each player is awarded a score based on the sum of values for the items they received. Otherwise, both of you get zero points.} \\
\texttt{(e) Your goal is to maximize the score you receive.} \\
\\ 
\texttt{Let us start.} \\
\\ 
\texttt{The set of available items is:} \\
\\ 
\texttt{1 bowl, 1 bottle, 2 fruits, 1 eraser} \\
\\ 
\texttt{Your secret value function is:} \\
\\ 
\texttt{bowl: 1, bottle: 0, fruit: 3, eraser: 3} \\
\\ 
\texttt{IMPORTANT: Your messages, unless it is the secret proposal, are directly transmitted to the other player, so do not include any response to the rules or text announcing your message. To make a secret proposal, use the indicated format. Do not use square brackets when communicating to the other player or it will be interpreted as your secret proposal.} \\
\\ 
\\ 
\texttt{<think>} \\
\texttt{Okay, let's see. I need to negotiate the division of these items: 1 bowl, 1 bottle, 2 fruits, 1 eraser. My value function is bowl:5, bottle:5, fruits and eraser:0. So I care a lot about getting the bowl and the bottle, but the fruits and eraser are worthless to me. The other player's value function is secret, so I don't know what they want. } \\
\\ 
\texttt{First, I should figure out what the other player might value. Since the items are bowl, bottle, fruits, and eraser, maybe they value the fruits and eraser? Or maybe they also want the bowl and bottle. But since there's only one of each, if both of us want them, we can't both get them. So the key is to find out what they want and see if we can split the items in a way that our proposals are complementary.} \\
\\ 
\texttt{The rules say that if our secret proposals are complementary, we get our scores. Otherwise, zero. So I need to make sure that whatever I propose, the other player's proposal doesn't conflict with mine. Since I don't know their value function, I have to make an educated guess.} \\
\\ 
\texttt{Let me think. The available items are limited. If I take the bowl and bottle, which are the only items I value, then the other player would have to take the fruits and eraser. But maybe they want some of the high{-}value items too. If they also want the bowl or bottle, then we might have a conflict. But since I can't get points from fruits or eraser, maybe I can offer them to the other player in exchange for the bowl and bottle. } \\
\\ 
\texttt{Wait, but how do I communicate that? I can send messages to negotiate. Let me start by sending a message to the other player. Maybe I can propose that I take the bowl and bottle, and they take the rest. But I need to phrase it in a way that encourages them to accept. } \\
\\ 
\texttt{So my first message could be: "I would like the bowl and the bottle. I don't value the fruits or eraser, so I'm happy for you to have them. Let me know if that works for you." That way, I'm stating my preference and offering the other items to them. If they don't need the bowl and bottle, they can take the fruits and eraser, and our proposals would be complementary} \\
            }
        }
    }
    & \\ \\

    \theutterance \stepcounter{utterance}  
    & & & \multicolumn{4}{p{0.6\linewidth}}{
        \cellcolor[rgb]{0.9,0.9,0.9}{
            \makecell[{{p{\linewidth}}}]{
                \texttt{\tiny{[GM$\langle$P2]}}
                \texttt{<think>} \\
\texttt{Okay, let's see. The user is trying to negotiate the division of items with another player. The available items are 1 bowl, 1 bottle, 2 fruits, and 1 eraser. The user's value function is bowl: 1, bottle: 0, fruit: 3, eraser: 3. Wait, no, wait. Wait, the user's secret value function is bowl: 1, bottle: 0, fruit: 3, eraser: 3. Wait, but in the initial problem statement, the user's value function was bowl:5, bottle:5, fruits and eraser:0. Wait, no, no. Wait, the user's secret value function is given as bowl: 1, bottle: 0, fruit: 3, eraser: 3. Wait, that's different from the initial problem statement. Wait, the user's value function is bowl: 1, bottle: 0, fruit: 3, eraser: 3. So the user values fruits and eraser highly, and the bowl a little, but the bottle is worthless. So the user's main goal is to get as many fruits and eraser as possible, but they can also take the bowl if needed.} \\
\\ 
\texttt{The other player's value function is secret. The user needs to negotiate with them to get a complementary proposal. The key is to figure out what the other player wants. Since the user's value function is known to them, but the other player's is not, the user needs to make a proposal that allows both to get their desired items without conflict.} \\
\\ 
\texttt{The available items are 1 bowl, 1 bottle, 2 fruits, 1 eraser. The user's value is 3 for each fruit and eraser, 1 for the bowl, 0 for the bottle. So the user's maximum possible score would be 3 (fruit) + 3 (fruit) + 3 (eraser) + 1 (bowl) = 10. But since there are only 2 fruits, maybe the user can get both fruits and the eraser, and maybe the bowl. But the other player might want the bowl or the bottle.} \\
\\ 
\texttt{The user's first message was to propose taking the bowl and bottle, but that's not optimal because the user doesn't value the bottle. Wait, the user's value for the bottle is 0} \\
            }
        }
    }
     \\ \\

    \theutterance \stepcounter{utterance}  
    & & & \multicolumn{2}{p{0.3\linewidth}}{
        \cellcolor[rgb]{0.9,0.9,0.9}{
            \makecell[{{p{\linewidth}}}]{
                \texttt{\tiny{[GM$|$GM]}}
                \texttt{continue} \\
            }
        }
    }
    & & \\ \\

    \theutterance \stepcounter{utterance}  
    & & & \multicolumn{2}{p{0.3\linewidth}}{
        \cellcolor[rgb]{0.9,0.9,0.9}{
            \makecell[{{p{\linewidth}}}]{
                \texttt{\tiny{[GM$|$GM]}}
                \texttt{<think>} \\
\texttt{Okay, let's see. The user is trying to negotiate the division of items with another player. The available items are 1 bowl, 1 bottle, 2 fruits, and 1 eraser. The user's value function is bowl: 1, bottle: 0, fruit: 3, eraser: 3. Wait, no, wait. Wait, the user's secret value function is bowl: 1, bottle: 0, fruit: 3, eraser: 3. Wait, but in the initial problem statement, the user's value function was bowl:5, bottle:5, fruits and eraser:0. Wait, no, no. Wait, the user's secret value function is given as bowl: 1, bottle: 0, fruit: 3, eraser: 3. Wait, that's different from the initial problem statement. Wait, the user's value function is bowl: 1, bottle: 0, fruit: 3, eraser: 3. So the user values fruits and eraser highly, and the bowl a little, but the bottle is worthless. So the user's main goal is to get as many fruits and eraser as possible, but they can also take the bowl if needed.} \\
\\ 
\texttt{The other player's value function is secret. The user needs to negotiate with them to get a complementary proposal. The key is to figure out what the other player wants. Since the user's value function is known to them, but the other player's is not, the user needs to make a proposal that allows both to get their desired items without conflict.} \\
\\ 
\texttt{The available items are 1 bowl, 1 bottle, 2 fruits, 1 eraser. The user's value is 3 for each fruit and eraser, 1 for the bowl, 0 for the bottle. So the user's maximum possible score would be 3 (fruit) + 3 (fruit) + 3 (eraser) + 1 (bowl) = 10. But since there are only 2 fruits, maybe the user can get both fruits and the eraser, and maybe the bowl. But the other player might want the bowl or the bottle.} \\
\\ 
\texttt{The user's first message was to propose taking the bowl and bottle, but that's not optimal because the user doesn't value the bottle. Wait, the user's value for the bottle is 0} \\
            }
        }
    }
    & & \\ \\

    \theutterance \stepcounter{utterance}  
    & & \multicolumn{4}{p{0.6\linewidth}}{
        \cellcolor[rgb]{0.9,0.9,0.9}{
            \makecell[{{p{\linewidth}}}]{
                \texttt{\tiny{[P1$\langle$GM]}}
                \texttt{<think>} \\
\texttt{Okay, let's see. The user is trying to negotiate the division of items with another player. The available items are 1 bowl, 1 bottle, 2 fruits, and 1 eraser. The user's value function is bowl: 1, bottle: 0, fruit: 3, eraser: 3. Wait, no, wait. Wait, the user's secret value function is bowl: 1, bottle: 0, fruit: 3, eraser: 3. Wait, but in the initial problem statement, the user's value function was bowl:5, bottle:5, fruits and eraser:0. Wait, no, no. Wait, the user's secret value function is given as bowl: 1, bottle: 0, fruit: 3, eraser: 3. Wait, that's different from the initial problem statement. Wait, the user's value function is bowl: 1, bottle: 0, fruit: 3, eraser: 3. So the user values fruits and eraser highly, and the bowl a little, but the bottle is worthless. So the user's main goal is to get as many fruits and eraser as possible, but they can also take the bowl if needed.} \\
\\ 
\texttt{The other player's value function is secret. The user needs to negotiate with them to get a complementary proposal. The key is to figure out what the other player wants. Since the user's value function is known to them, but the other player's is not, the user needs to make a proposal that allows both to get their desired items without conflict.} \\
\\ 
\texttt{The available items are 1 bowl, 1 bottle, 2 fruits, 1 eraser. The user's value is 3 for each fruit and eraser, 1 for the bowl, 0 for the bottle. So the user's maximum possible score would be 3 (fruit) + 3 (fruit) + 3 (eraser) + 1 (bowl) = 10. But since there are only 2 fruits, maybe the user can get both fruits and the eraser, and maybe the bowl. But the other player might want the bowl or the bottle.} \\
\\ 
\texttt{The user's first message was to propose taking the bowl and bottle, but that's not optimal because the user doesn't value the bottle. Wait, the user's value for the bottle is 0} \\
            }
        }
    }
    & \\ \\

    \theutterance \stepcounter{utterance}  
    & \multicolumn{4}{p{0.6\linewidth}}{
        \cellcolor[rgb]{0.9,0.9,0.9}{
            \makecell[{{p{\linewidth}}}]{
                \texttt{\tiny{[P1$\rangle$GM]}}
                \texttt{<think>} \\
\texttt{Okay, let me try to figure this out. The user's value function is bowl:1, bottle:0, fruit:3, eraser:3. So they really want the fruits and eraser, and the bowl is a minor point. The other player's value function is secret, but the user needs to make a proposal that's complementary. The key is to figure out what the other player might want.} \\
\\ 
\texttt{The available items are 1 bowl, 1 bottle, 2 fruits, 1 eraser. The user's main goal is to get as many fruits and eraser as possible. Since there are two fruits, maybe the user can take both, plus the eraser. That would give them 3+3+3=9 points. The bowl is worth 1, but the user might not need it unless the other player wants it. The bottle is worthless to the user, so they can let the other player have it.} \\
\\ 
\texttt{But the other player's value function is unknown. If the other player values the bowl or bottle, the user needs to make sure their proposal doesn't conflict. For example, if the other player wants the bowl, the user could take the bowl, but since it's only worth 1, maybe it's better to let the other player have it if they need it. The user's main assets are the fruits and eraser.} \\
\\ 
\texttt{So the user's optimal strategy is to take both fruits and the eraser, which gives them 9 points. The bowl and bottle can be given to the other player. But the user needs to check if the other player's proposal would conflict. If the other player takes the bowl and bottle, that's fine because the user doesn't value them. So the user's proposal should be {[}Proposal: 2 fruit, 1 eraser{]}. That way, they get the high{-}value items, and the other player can take the bowl and bottle. If the other player's proposal is for the bowl and bottle, then the proposals are complementary, and both get their items. If the other player also wants fruits or eraser, there might be a conflict, but since the user is taking all the fruits and eraser, the other player can't get those. But the user's main goal is to maximize their own score, so taking all the high{-}value items is the way to go.} \\
\texttt{</think>} \\
\\ 
\texttt{{[}Proposal: 2 fruit, 1 eraser{]}} \\
            }
        }
    }
    & & \\ \\

    \theutterance \stepcounter{utterance}  
    & & & \multicolumn{2}{p{0.3\linewidth}}{
        \cellcolor[rgb]{0.9,0.9,0.9}{
            \makecell[{{p{\linewidth}}}]{
                \texttt{\tiny{[GM$|$GM]}}
                \texttt{proposal} \\
            }
        }
    }
    & & \\ \\

    \theutterance \stepcounter{utterance}  
    & & & \multicolumn{2}{p{0.3\linewidth}}{
        \cellcolor[rgb]{0.9,0.9,0.9}{
            \makecell[{{p{\linewidth}}}]{
                \texttt{\tiny{[GM$|$GM]}}
                [0, 0, 2, 1]
            }
        }
    }
    & & \\ \\

    \theutterance \stepcounter{utterance}  
    & & \multicolumn{4}{p{0.6\linewidth}}{
        \cellcolor[rgb]{0.9,0.9,0.9}{
            \makecell[{{p{\linewidth}}}]{
                \texttt{\tiny{[GM$\rangle$P2]}}
                \texttt{The other player has submitted a secret proposal. It is now your turn to submit a proposal of your own.} \\
            }
        }
    }
    & \\ \\

    \theutterance \stepcounter{utterance}  
    & & & \multicolumn{4}{p{0.6\linewidth}}{
        \cellcolor[rgb]{0.9,0.9,0.9}{
            \makecell[{{p{\linewidth}}}]{
                \texttt{\tiny{[GM$\langle$P2]}}
                \texttt{<think>} \\
\texttt{Okay, let's see. The other player has already submitted a secret proposal, so I need to come up with my own. The available items are 1 bowl, 1 bottle, 2 fruits, and 1 eraser. My value function is bowl:1, bottle:0, fruit:3, eraser:3. So I want to maximize my score by getting as many fruits and erasers as possible, since they give me 3 each. The bowl is worth 1, and the bottle is worthless.} \\
\\ 
\texttt{The other player's proposal isn't visible to me, but I need to make sure that our proposals are complementary. If they asked for some items, I need to take the remaining ones. Since I don't know their proposal, I have to make an educated guess. } \\
\\ 
\texttt{If I take both fruits and the eraser, that's 3*2 + 3 = 9 points. The bowl is worth 1, so maybe I can take that as well. But there's only one bowl. If the other player also wants the bowl, then we can't both have it. But since the bottle is worth 0 to me, maybe I can let them have the bottle if they don't need the bowl. } \\
\\ 
\texttt{Wait, the other player's value function is secret. They might value the bottle, but since I don't care about it, I can offer it to them. So my proposal should include the items I value most, and leave the ones I don't need for them. } \\
\\ 
\texttt{So my best bet is to take both fruits and the eraser, and maybe the bowl. Let me check the quantities. There are 2 fruits, 1 eraser, 1 bowl. If I take all of those, that's 2 fruits, 1 eraser, 1 bowl. The remaining item is the bottle. Since the other player can take the bottle, our proposals would be complementary. } \\
\\ 
\texttt{So my proposal would be {[}Proposal: 2 fruit, 1 eraser, 1 bowl{]}. That way, I get 2*3 + 3 + 1 = 10 points. The other player gets the bottle, which they might value. If their proposal doesn't include the bottle, then we're good. If they also wanted the bowl, but I took it, then their proposal might conflict. But since I don't know their values, I have to assume they might not} \\
            }
        }
    }
     \\ \\

    \theutterance \stepcounter{utterance}  
    & & & \multicolumn{2}{p{0.3\linewidth}}{
        \cellcolor[rgb]{0.9,0.9,0.9}{
            \makecell[{{p{\linewidth}}}]{
                \texttt{\tiny{[GM$|$GM]}}
                \texttt{proposal} \\
            }
        }
    }
    & & \\ \\

    \theutterance \stepcounter{utterance}  
    & & & \multicolumn{2}{p{0.3\linewidth}}{
        \cellcolor[rgb]{0.9,0.9,0.9}{
            \makecell[{{p{\linewidth}}}]{
                \texttt{\tiny{[GM$|$GM]}}
                [1, 0, 2, 1]
            }
        }
    }
    & & \\ \\

    \theutterance \stepcounter{utterance}  
    & & & \multicolumn{2}{p{0.3\linewidth}}{
        \cellcolor[rgb]{0.9,0.9,0.9}{
            \makecell[{{p{\linewidth}}}]{
                \texttt{\tiny{[GM$|$GM]}}
                [[0, 0, 2, 1], [1, 0, 2, 1]]
            }
        }
    }
    & & \\ \\

\end{supertabular}
}

\end{document}
